\documentclass[12pt]{article}

\def\version{1.2}

\usepackage {manual}
\usepackage {fullpage}
\usepackage {makeidx}

\expandafter \ifx \csname driver\endcsname\relax \def \driver{}\fi 
\usepackage [\driver]{hyperref}

\usepackage {graphicx}
\usepackage {advi}
\usepackage {advi-annot}
\usepackage {bubble}

% For splash inclusion
\usepackage{color}
\usepackage{graphicx}
\usepackage{tabularx}
\usepackage{pst-all}
\usepackage{hyperref}
% End of splash specific stuf

\newdimen \tmpdim
\title {\setbox0 \hbox {\bf\ActiveDVI}\global \tmpdim \wd0
\advance \tmpdim by 4em
\begin{tabular}{c}
\box0 \\
\includegraphics[width=\the\tmpdim]{../tex/bar.jpg.eps}\\
Reference manual\\
Version {\version}
\end{tabular}
}
\author {Didier R{\'{e}}my and Pierre Weis}

\thispagestyle{empty}

\begin{document}

\maketitle

\def \abstractname {}
\begin{abstract}
{\bf\ActiveDVI} is a viewer for DVI files that also recognizes a new class of 
\docdef \special's targeted to presentations via laptop computers:
various special effects can easily be incorporated to the presentation,
via a companion \doctt{advi} {\LaTeX} package. 

\bigskip
{\ActiveDVI} is copyright \copyright\space 2001, 2002 INRIA
and distributed under the Gnu Library General Public License ---see the LGPL
file included in the distribution.
\end{abstract}


\vfill
 \section* {Acknowledgements and contributors}

{\ActiveDVI} is based on Mldvi, written by Alexandre~Miquel, which
constitutes its core. Mldvi alone is distributed 
under LGPL license and is available at \url{http://pauillac.inria.fr/~miquel/}. 

{\ActiveDVI} has been developed by Jun~Furuse, Didier~R{\'{e}}my and
Pierre~Weis with also contributions by Roberto~Di~Cosmo, Didier~Le~Botlan,
Xavier~Leroy, and Alan~Schmitt.

\clearpage
\advitoc
\tableofcontents
\endadvitoc
\clearpage

\section {Installation}

\section {Safety concerns when using the {\ActiveDVI} previewer}

\paragraph {\bf Warning!}

{\ActiveDVI} may execute programs and commands embedded into the DVI file.
Hence, when playing a DVI file from an untrusted source,
you should run \doctt{advi} with the \doctt{-safer} option, 
that inhibits the execution of embedded applications. 

This warning applies in particular if you choose {\ActiveDVI} as your default 
{\tt meta-mail} previewer for the {\tt application/x-dvi} mime-type.

Conversely, you may run \doctt{advi} with the \doctt{-exec} option to
automatically launch embedded applications.

The last safety option is the default one: the \doctt{-ask} option
tells {\ActiveDVI} to ask the user each time it has to launch an
application. (Note that in such a case {\ActiveDVI} asks only once to
launch a given application: it remembers your previous decisions
concerning a command and acts accordingly for the rest of the
presentation.) 

The safe \doctt{-ask} option is set by default. It can also be set for
all the users by the machine administrator or on a local scale by
individual users, and even via a per directory setting ! (Conveninient
to gracefully run your talks, while still be cautious when running
talks from others.)

\section {Using the {\ActiveDVI} previewer}

\subsection {Command line options}

{\ActiveDVI} is invoked with the following command syntax
\begin{verbatim}
advi [options] dvifile
\end{verbatim}

\noindent The \verb"advi" commands recognized the following options:

\medskip\noindent\begin{tabular}{ll}
\verb"-help"               & Short command line options help
\end{tabular}

\subsubsection*{Window and display specifications}

\medskip\noindent\begin{tabular}{ll}
\verb"-geometry geom"      & Geometry of \ActiveDVI's window specification
\end{tabular}

Geometry \verb"geom" is specified in pixels, using the standard format
for XWindow geometry specifications (i.e:
\verb"WIDTHxHEIGHT[+XOFFSET+YOFFSET]").

\medskip\noindent\begin{tabular}{ll}
\verb"-fullwidth"          & Adjust the size of the window to the width of the screen
\end{tabular}

\medskip\noindent\begin{tabular}{ll}
\verb"-nomargins"          & Cancel horizontal and vertical margins \\
\verb"-hmargin dimen"      & Horizontal margin specification (default: 1cm) \\
\verb"-vmargin dimen"      & Vertical margin specification   (default: 1cm) \\
\end{tabular}

Dimensions are specified as numbers optionally followed by two letters
representing units. When no units are given, dimensions are treated
as numbers of pixels. Currently supported units are the standard TeX
units as specified in the TeXbook (D. Knuth, Addison-Wesley, (C)
1986):
 `pt' (point), `pc' (pica), `in' (inch), `bp' (big point),
 `cm' (centimeter), `mm' (millimeter), `dd' (didot point),
 `cc' (cicero) and `sp' (scaled point).
Note that dimensions are specified w.r.t the original TeX document,
and do not correspond to what is actually shown on the screen, which
can be displayed at a different resolution than specified in the
original TeX source.

\medskip\noindent\begin{tabular}{ll}
\verb"-crop"               & Crop the window to the best size (default) \\
\verb"-nocrop"             & Disable cropping
\end{tabular}

\subsubsection*{Helpers specification}

\medskip\noindent\begin{tabular}{ll}
\verb"-pager" & Specify the name of the pager to launch on a \verb"txt" link\\
\verb"-browser" & Specify the name of the browser to a \verb"html" link\\
\end{tabular}

\subsubsection*{Debugging options}

\medskip\noindent\begin{tabular}{ll}
\verb"--debug" & \\
\verb"--debug_pages" & \\
\verb"--show_ps" & \\
\verb"--verbose_image_access" & \\
\end{tabular}

\subsubsection*{Rendering options}

\medskip\noindent\begin{tabular}{ll}
\verb"-A" & Toggle Postscript antialiasing \\
\verb"-passive" & Inhibit effects that are visible when redrawing the page
\\ & (Transitions, delays, embeded apps) \\
\end{tabular}

\subsubsection*{Safety options}
\medskip\noindent\begin{tabular}{ll}
\verb"-exec" & Set automatic execution safety policy \\
\verb"-ask" & Set automatic execution safety policy \\
\verb"-safer" & Set no execution of embedded aplications safety policy \\
\end{tabular}

\subsubsection*{Option files option}
\medskip\noindent\begin{tabular}{ll}
\verb"-option-file <filename>" & Load {\tt filename} as a file containing a list of options \\
& as given on the command line to \doctt{advi} \\
\end{tabular}


\subsection {Cut and paste}

Text can be copied from the {\ActiveDVI} previewer to an another application.
However, this uses the \doctt{XBuffer} and not the \doctt{XSelection}
mechanism.
\begin {itemize}
\item
Shift middle-click copies the current word. 
\item
Shift right-click and drag copies the specified region. 
\end {itemize}
Moreover, Shift left-click dump an ASCII representation of click 
in the source file. This expects the DVI to be instrumented with line
numbers of the form
\begin{quote}
\doctt {line:} \docid {line} \docid {file}
\end{quote}
where \docid {line} and \docid {file} are the current source line and
current source file.

The position is exported in ASCII in the form
\begin{quote}
\verb"#line "\docid{before}, \docid{after} 
\doctt {<<}\docid{prefix}\doctt{>><<}\docid{suffix}\doctt{>>}
\docid{file} 
\end{quote}
where \docid {before} and  \docid {after} are the enclosing line numbers, 
\docid{prefix}\docid{suffix} are the word constituent surrounding the
position, and  file is the current file. 

Line numbers default to $0$ when not found. Note that line numbers may be
inconsistent if they \docdef \special-line commands have not been inserted
close enough to the position.

\subsection {Hyperref}

{\ActiveDVI} supports {\LaTeX} {\tt hyperref} package 
with both internal and cross-file references. For cross-file
references, it launches a new \verb"advi" process. 


{\ActiveDVI} improves the treatment of hyperrefs over convential previewers,
by emphasizing the hypertarget text of an hyperlink. Thus, an hypertarget
definition \docdef\hypertarget \docarg{tag}\-\docarg{text}
should make the activation of the link \docid{tag}
not only move to the page where \docid{text} occurs, but also emphasize
\docid{text}. 
%
However, since {\docdef\hypertarget} does include its second argument
withing the target, we use the following command instead:
\begin{verbatim}
\newcommand{\softtarget}[2]%
  {\special{html:<A name=\hyper@quote #1\hyperquote>}#2\special{html:</A>}}
\end{verbatim}
(If you are viewing this document with \softtarget{a}{\ActiveDVI}, you may
move over or click on \hyperlink{a}{this area} to see the effect.) 

\subsection {Floating table of content and thumbnails}

There are two ways to include a floating table of content while previewing. 
\begin {itemize}

\item
{\ActiveDVI} recognizes the reserved hypertargets \doctt{/toc.first}
and \doctt{/toc.last} as markers for the first and last page of the table
of contents, which will become floating pages, {\em i.e.} accessible from
anywhere in the document with the default keybinding \verb't'.
The first stroke on \verb't' will show the first page of the table of
content. Successive stokes will show the following pages. Prefix argument 
may also be used to directly access second of third pages. 

The package \softlink{advi}{{\tt advi}} below redefines
{\docdef \tableofcontents} so that it automatically inserts the above
markers around the table of content. It also provides two macros, {\docdef
\advitoc} and \docdef \endadvitoc, that can be used to insert these markers
explicitly when the table of context is hand-made.

\item
If no table of content markers are found, then {\ActiveDVI} will
compute thumbnails, {\em i.e.} will present pack several pages at a smaller 
scale on a single page. {\ActiveDVI} will calculate the scale needed to 
fit all thumbnails on a single page, and will use it if it is 
less or equal to a maximal value, or take the maximal value otherwise
and use several pages. The default maximal scale value is 5, so that
25 thumbnails can fit on the same page. This value can be changed with the
option \doctt{-thumbnail\_scale}.  

By default, thumbnails are made for all pages. However, thumbnail pages can
be defined explicitly with an hypertarget whose anchor is of the form
\doctt{/page.\docid{suffix}}. In this case, all thumbnail must be explicitly
given.

\end {itemize}
By default, the binding \doctt{T} processes thumbnails and the binding
\doctt{t}  displays thumbnails if already processed or show the table of
content if available, otherwise does nothing. Thumbnails must be made
explicitlty, to prevent accidental hit of the key, which may then 
take some time.


\subsection {Moving around}

See the key bindings in the appendix.

\subsection {Using and making special effects}

Presentation examples can be found in the 
\verb"examples" directory. Don't miss to play them! Then, feel free
to read their source code and copy the effects they provide.

{\ActiveDVI} can be used as is, but will shine when driven by a user
with a bias toward programming: special effects can easily 
be realized by using the \LaTeX package provided with the distribution. 

Creative advanced users may program the presenter at various levels, either
using or defining simple {\LaTeX} macros, writing new {\LaTeX} package
files, or by implementing extension to the previewer itself.

% If you ever succeed at adding a nice special effect that could be of
% a general interest, or if you can provide a nice and fancy presentation
% of your own that can serve as an example to the \ActiveDVI's users
% community, please don't hesitate to drop a mail to the implementors
% at \verb"advi@pauillac.inria.fr".


\section {The \softtarget{advi}{{\tt advi.sty}} \LaTeX~package}


{\ActiveDVI} provides some {\LaTeX} packages to facilitate animations and
interaction with the presenter from within your {\LaTeX} source text.

The  advi package is the main package to include when writing
a presentation for {\ActiveDVI}. It provides the main set of interactive
commands for {\ActiveDVI} to animate the show. However, there is no
need to load the package if no {\ActiveDVI} special effects are
required for the presentation.

\subsection{Printing the presentation}

The {\tt advi} package recognizes the special option \verb"ignore",
which is devoted to help the production of a printable version of the
presentation: the \verb"ignore" option makes the package not to
produce {\tt advi} specials, so that the show can be previewed by
other previewers or turned into Postscript with {\tt dvips}. Of
course, this option disables most effects that cannot be printed,
although some of them are still approximated.

If the {\tt ignore} option is not set globally, it can be set
locally with the commands {\docdef \adviignore}. However, this will not
prevent all effects, since some decisions are taken when the package is
loaded.

The package also defines the conditional {\docdef \ifadvi} which
evaluates its first argument if advi is not in ignore mode and its
second argument otherwise.

\subsection {Pause, Record and Play}

\medskip\noindent
\docdef \adviwait\docopt{seconds}
\begin{quote}
Wait for \docid{seconds}. If no argument is provided, waits until the
user requests to continue (hitting a key to move to next pause or to
change page).
\end{quote}

\medskip\noindent
\docdef \advirecord \doctt{[play]}\docarg{tag}\docarg{latex code}
\\
\docdef \begin\doctt{\{advirecording\}}%
\doctt{[play]}\docarg{tag}\docarg{latex code}%
\docid {text}\docdef \end\doctt{\{advirecording\}}
\begin{quote}
Processes \docid{latex code} and records the corresponding DVI output,
bound to the tag \docid{tag}. When recording, the DVI output is
not displayed unless the option \doctt{play} is set.

Records can be nested. If so, the inner record is always recorded
with its own  tag; if the inner record is played when recording, it is
also recorded as part of the outer tag.

If the environment form is used, the \docid{latex code} may contain
fragile commands.
\end{quote}

\medskip\noindent
\docdef \adviplay\docopt{color}\docarg{tag}
\begin{quote}
Replay the DVI previously recorded and bound to \docid{tag}.
The optional argument changes the color to \docid{color} during replay. 
\end{quote}

\medskip\noindent
\docdef \advianchor\docopt{activation}\docarg{tag}\docarg{text}
\\
\docdef \begin\doctt{\{advianchoring\}}\docopt{activation}\docarg{tag}%
\docid {text}\docdef \end\doctt{\{advianchoring\}}

\begin{quote}
Plays the record bound to \docid{tag} when the anchor is activated.

The argument \docid{activation} may either be \doctt{over} or
\doctt{click}. The page is reset to its original appearance when the
anchor is no more activated (the mouse leaves the anchor area or the
button is released).

If the environment form is used, \docid{text} may contain fragile commands. 

\end{quote}


\subsection {Images}

Images can be encapsulated into the presentation using the package
CamlImages provided with the distribution
(see section~\ref {advi-grahicx})

\subsection {Background}

Background of pages can be set in the {\LaTeX} source of the page,
either as a plain color or as an image (or both!).
You can specify a global option to the background setting, so that this
background will be used for the remaining pages of the presentation
(otherwise the presenter resets background options at each new page).

Background images can be lighten by specifying an alpha value (a
floating point number between 0 and 1) that mesures the mixing
between the color of the background and the image.

Background images can also be blended, meaning that you can choose the
algorithm that surimposes the image to the background.

\medskip\noindent
\docdef \advibg \doctt {[global]}\docarg{decl}
\begin{quote}
where \docid{decl} is a list of settings of the following from:

\dockey[none]{color}{color}
\begin{quote}
Set the background color to \docid{color}.
If \docid{color} is \doctt{none} this unsets the background color. 
Otherwise, \docid{color} must follow the convention of the package 
\doctt{color}, that is, it should either be a previously declared color or 
of the form \docopt{model}\docarg{model color specification}. 

For example, the following specifications are all correct:
\begin{quote}\tt
color=blue\\
color=[named]\{Yellow\}\\
color=[rgb]\{0.7,0.3,0.8\}\\
\end{quote}
\end{quote}

\dockey[none]{image}{file}
\begin{quote}
Use the image found in \docid{file} as background (\doctt{none} means unset).
\end{quote}

\dockey[auto]{fit}{fit style}
\begin{quote}
Fit the background image according to \docid{style}, which may be 
one of the following keywords:
\begin{center}
\tt
auto \hfill {\rm or} \hfill
\begin{tabular}{l@{\qquad}c@{\qquad}r}
topleft&   top& topright\\
left&  center&  right \\
bottomleft& bottom & bottomright\\
\end{tabular}
\end{center}
The {\tt auto} fit style means scaling the image as desired in both directions
so that it fits the entire page. Other styles only force the same scaling
factor in both directions:
\begin {itemize}

\item Corner-styles means set the image in the corresponding corner
and scale it to cover the entire page.

\item {\tt center} means set the image in the center of the page and
scale it to cover the entire page.

\item Segment-styles means adjust the image and the page on the
segment (in which case, the image may not completely cover the page on
the opposite side).

\end {itemize}
\end{quote}

\dockey[none]{alpha}{float}
\begin{quote}
Set the alpha channel factor for the background image to \docid{float}
(\doctt{none} means unset). An alpha factor of $0$ means that the
image is not visible at all; conversely, an alpha factor of $1$ means that the
image covers the background.
\end{quote}

\dockey[none]{blend}{blend mode}
\begin{quote}
Set the blend mode to \docid{blend mode}, which should be one of the
following (unchecked): {\tt difference}, XXX?
(\doctt{none} means unset).
\end{quote}

\goodbreak
\dockey{none}{}
\begin{quote}
Unset all background parameters.
This key must appear on its own, no arguments or keys are allowed. 
\end{quote}
The optional parameter \doctt{global} indicates that
the definition is global and will affect the following pages, 
as well as the current page.

By default, the background settings only affect the current page.
\end{quote}

\subsection {Transitions}

\docdef \advitransition \doctt {[global]}\docarg {decl}
\begin{quote}
where \docid{decl} is a list of settings of the following from:

\dockey{none}{} or \dockey{slide}{} or \dockey{block}{} or \dockey{wipe}{}
\begin{quote}
Set the transition mode to the corresponding key.
One of this key is mandatory (several is should be illegal, but in fact
overrides the previous one). 
\end{quote}

\dockey{from}{direction}
\begin{quote}
Make the transition come from \docid{direction}.  Directions are not checked
but should be one of the following: 
\begin{center}
\tt
\begin{tabular}{l@{\qquad}c@{\qquad}r}
topleft&   top& topright\\
left&  center&  right \\
bottomleft& bottom & bottomright\\
\end{tabular}
\end{center}
\end{quote}
The default direction, to be used when no local or global direction has been
specified, is determined dynamically: \verb"right" when coming from
previsous page, \verb"left" when coming from next page, and \verb"top"
otherwise.

\dockey{steps}{n}
\begin{quote}
Make the transition in \docid{n} steps. 
\end{quote}
As for \docdef \advibg, the optional parameter \doctt{global} indicates that
the definition is global and will affect the following pages, 
as well as the current page. 

By default, the transition definitions only affect the current page. 
\end{quote}

\medskip\noindent
\docdef \advitransbox \docarg {key=val list}\docarg{hbox material}
\begin{quote}
where \docid {key=val list} is as above and \docarg{hbox material} is
whatever can follow an \docdef \hbox{} command. In particular, the material
may contain verbatim commands, since as for the \docdef\hbox{} it is 
parsed incrementally.
\end{quote}

The optional parameter \doctt{global} indicates that
the definition is global and will affect the following pages, 
as well as the current page.

By default, the transition affects only the current page.

\subsection {Embedded applications}

{\ActiveDVI} can launch arbitrary applications you need to animate the
show.

The {\LaTeX} command is 

\medskip\noindent
\docdef \adviembed\docopt{key=value list}\docarg{command} 

\begin{quote}
where \docid{key=value list} is a list of bindings of the following kind:

\dockey{name}{name}
\begin{quote}
Allows to refer to the embedded application as \docid{name}. 
Anonymous applications have actually the default name \doctt{anonymous}.
\end{quote}

\dockey{ephemeral}{name}

\begin{quote}
This is the default case: the embedded application is killed when the page
is turned.
\end{quote}

\dockey{persistent}{name}

\begin{quote}
The application keeps running in the background, but will only be visible in
the page where it has been launched.
\end{quote}

\dockey{sticky}{name}
\begin{quote}
The application keeps running and remains visible when turning pages.
\end{quote}

\dockey{width}{dim}\\
\dockey{height}{dim}
\begin{quote}
The application takes \docid{dim} width (respectively height) space in 
latex. Both values default to 0pt.

These dimensions are also substituted for all occurrences of \verb"!g!" in
the command string. 
\end{quote}

\end{quote}

\medskip\noindent
\docdef \advikillembed \docarg{name}
\begin{quote}
Kill all embedded applications named \docid{name}.
An optional signal value or symbolic name can be given to send to the
designed process: for instance \verb"\advikillembed[SIGUSR1]{clock}"
will send the \verb"SIGUSR1" signal to the embedded application named
\verb"clock".

Signal value defaults to -9.
\end{quote}


\subsection {Active anchors}

The command is 

\medskip\noindent
\docdef \advianchor \docopt{decl}\docarg{tag}\docarg {text}
\\
\docdef \begin\doctt{\{advianchor\}}\docopt{decl}\docarg{tag}%
\docid {text}\docdef \end\doctt{\{advianchor\}}
\begin{quote}
The text is first display as usual, then the anchor is drawn
according to the style given by \docid{decl}, and made active. 
It's acctivation, which depends on the mode give by \docid{decl}, 
will play record \docid{tag}. 

The declarations \docid{decl} are of the following form:
\begin{description}
\item [\doctt{over} \textrm{or} \doctt{click}]\indent\par 
this defines the mode of  activation, which is either by
moving mouse  over  or clicking on the anchor. 
The default mode is \doctt{over}.
\item [\doctt{box}, \doctt{invisible}, or \doctt{underline}]\indent\par
the defines the style in which the anchor should be drawn.
The default style is \doctt{boxed}.
\end{description}
In the environment form, \docid{text} may contain fragile commands. 
\end{quote}

\medskip\noindent
\docdef \adviemphasize \docopt{color}\docarg {text}
\begin{quote}
This makes an invisible anchor around \docid{text}, which when activated
will redraw text in a box colored with \docid{color}, which defaults to
\doctt{yellow}. 
\end{quote}

\subsection {Postscript specials}

{\ActiveDVI} can deal with most of PStricks calling {\tt ghostscript} on
included Postscripts. However, the interactive between {\ActiveDVI} and
Postscripts is not always working.
\begin {itemize}

\item
Since characters are rendered by {\ActiveDVI}, some PStricks are not allowed.
\item
Some change of repairs are also not yet correctly performed.  
\end {itemize}

\subsubsection {Overlays}

\index{overlays}

The {\tt overlay} class implements overlays with PStricks.  Instead,
{\ActiveDVI} implements overlays directly, using records and plays.
This is more efficient, and of course more natural.
(In fact, {\ActiveDVI} chooses the cumulative semantics of overlays
display if all layers below the current overlay.)

The {\tt xprosper} style derived from the {\tt prosper} class uses 
the {\tt overlay} class and also work with {\ActiveDVI} in exactly the same
way (relaxing the \overlay@loop macro inhibits all layers, but the first
page). 

\subsubsection {PStricks}

{\ActiveDVI} can deal with most of PStricks. 
\begin {itemize}

\item[+]
Simple drawings work

\item[+] 
{\tt\string\SpecialCoor} works, {\em i.e.} commands of the form
{\tt \string \rput \{A\}\{bla bla\}} works where {\tt A} is a node

\item[+]
Connections between nodes {\tt \string \ncarc}, {\tt \string \ncarc},
also works.
\end {itemize}
However, some PStricks are known not to work.
\begin{itemize}
\item[--]
Labels over arrows \verb"\Aput", \verb"Bput", etc.
(they change the Postscript coordinates...)

\item[--]
{\tt pspicture}
(idem, but drawings that are not embedded in pspictures works)

\end {itemize}


\section {Auxiliary \LaTeX~packages}

\subsection{The superpose package}

This package allows to easily superpose horizontal material, creating 
the smallest horizontal box to fit all of the superpositions.
\begin{verbatim}
\usepackage{superpose}
\end{verbatim}
The package defines a single environment:

\medskip\noindent
\docdef\begin\doctt{\{superpose\}}\docopt{alignment}\docid {list}%
\docdef\end\doctt{\{superpose\}}
\begin{quote}
The \docid{alignment} can be one the letters \doctt c (default value), \doctt
l, or \doctt r.

Items of the list are separated by \verb"\\" as in tabular environments.
Each item should be a horizontal material.
\end{quote}

\subsection{The bubble package}

This package allows to easily draw bubbles over some text. 
\begin{verbatim}
\usepackage{bubble}
\usepackage[ps]{bubble}
\end{verbatim}
By default bubbles are produced using the \doctt{epic} and \doctt{eepic}
packages, for portability. However, for better rendering and easier 
parameterization, bubbles can also be drawn using the \doctt{ pst-node}
package of the pstricks collection. This is what the \doctt{ps} option
is designed for.

The package defines a single command:

\medskip\noindent
\docdef \bubble \docopt{key=value list}\docarg{anchor}%
    \docopt{ps options}\docpar {pos}\docarg{text}
\begin{quote}
The \docid{key=value list} is a list of bindings of the following kind:

\dockey[yellow]{bg}{color}
\begin{quote}
The background color for annotations. 
\end{quote}

\dockey[yellow]{unit}{dim}
\begin{quote}
Set the package unit to \docid{dim}.
\end{quote}

\dockey[c]{col}{colspec}
\begin{quote}
Where \docid{colspec} is a column specification for the tabular
environment. Moreover, the following abbreviations are recognized:

\smallskip
\tt
\begin{tabular}{|l|l|}
\hline
\rm\em key & \rm expands to \\
\hline
c &  col=c\\
l &  col=l\\
r &  col=r\\
p=\docid{w} & col=p\docarg{w}\\
\cline{1-2}
\end{tabular}
\begin{tabular}{|l|l|}
\hline
\rm\em key & \rm expands to \\
\hline
C &  col=\{>\{\$\}c<\{\$\}\}\\
L &  col=\{>\{\$\}l<\{\$\}\}\\
R &  col=\{>\{\$\}r<\{\$\}\}\\
P\docid{w} &  \doctt{col=\{>\{\$\}p\docarg{w}<\{\$\}\}}\\
\cline{1-2}
\end{tabular}
\end{quote}

\docid{pos} is the optional relative position of the annotation, 
it defaults to $1,1$, and is counted in the package units. 

\docid{ps options} are passed to the command \docdef\psset) 
in \doctt{ps} mode and ignored otherwise. 
\end{quote}
Parameters (color and tabular columns specifications) can also 
be set globally using the command:

\medskip \noindent
\docdef \setkeys\doctt{\{bubbleset\}}\docarg{key=value list}

\subsection{The annotations package}

This package uses active anchors and the bubbles package to 
allow annotations by raising a bubble when the cursor is over the anchor.

The package provides a single command

\medskip\noindent
\docdef \adviannot \docopt{key=value list}\docarg{anchor}%
    \docopt{ps options}\docpar {pos}\docarg{text}
\begin{quote}
whose options are identical to those of \docdef \bubble, but the bubble
appears withing an active anchor. 
\end{quote}

%Annotation location can also be invisible \advi{C}{\tt{ }}instead of visible.

\subsection{\label {advi-grahicx}The advi-graphicx package}

This 3-line package provides loads the graphicx prackage and provides
declarations so that JPEG, EPS, TIF, TIFF source images can be embedded:
{\ActiveDVI} will preview these directly while other drivers will translate
them on demand. 

\appendix

\section {Limitations}

\subsection*{Postscript Fonts}

Postscript fonts are not currently supported. 

\subsection*{Inlined Postscript and Ghostscript}

PS relies on {\tt ghostscript} to display Postscript inlined specials.
However, some earlier releases of {\tt ghostscript} 
implements Postscript {\tt flushpage} command as a {\tt XFlush} call
which does not force the evaluation of commands, and thus makes the 
synchronization between {\tt ghostscript} and {\ActiveDVI} drawings 
uncontrollable. In this case, the interleaving of inlined postscript 
and other material may be inconsistent. 

Fortunately, recent versions of ghostscript $(> 6.5)$ have
fixed this problem by using {\tt XSync(false)} instead.
If you use those versions of ghostscript, inlined specials should be
correctly rendered.

Unfortunately, some releases of version 6.5x also carry a small but
fatal bug for \ActiveDVI, that will hopefully be fixed in future
releases. A workaround is available here
\url{http://cristal.inria.fr/~remy/ghostscript/}.

\subsection*{Inlined Postscript change of coordinates}

So far, the implementation of inlined Postscript does not correctly handle
complex change of coordinates.  (See PStricks section). 

\section {Bug Reports}

Please, send bug reports to 
\url{mailto:advi@pauillac.inria.fr}. 
\\
See \url{http://pauillac.inria.fr/advi} for up to date information.


\section {Key bindings}

% Macros specific to splash.tex
\definecolor{red}{rgb}{1.0,0.0,0.0}
\def\key#1{\textcolor{red}{#1}}
\def\ikey#1#2#3{\key{#1} & #2 & -- & #3\\}
\definecolor{blue}{rgb}{0.0,0.0,1.0}
\def\arg{\textcolor{blue}{arg }}

\let \Newpage \newpage
\def \newpage {}

Advi recognizes the keystrokes listed below when typed in its window.
Some keystrokes may optionally be preceded by a number, called \arg
below, whose interpretation is keystroke dependant. If \arg is unset,
its value is 1, unless specified otherwise.

Advi maintains an history of previously visited pages organized as a stack. 
Additionnally, the history contains marked pages which are stronger than
unmarked pages. 

\vspace{\stretch{1}}

\newpage

\vspace{\stretch{1}}

\subsection*{Survival command kit}

\noindent
\begin{tabularx}{\linewidth}{clcX}
\ikey{?}{info}{This quick info and key bindings help.}
\ikey{q}{quit}{End of show.}
\ikey{space}{continue}
{Move forward (\arg pauses forward if any, or do as \key{return} otherwise).}
\end{tabularx}

\vspace{\stretch{1}}

\subsection*{Moving to pages}

\noindent
\begin{tabularx}{\linewidth}{clcX}
\ikey{n}{next}
{Move \arg physical pages forward, leaving the history unchanged.}
\ikey{p}{previous}
{Move \arg physical pages backward, leaving the history unchanged.}
\ikey{,}{begin}{Move to the first page.}
\ikey{.}{end}{Move to the last page.}
\ikey{g}{go}
{If \arg is unset move to the last page.
 If \arg is the current page do nothing.
 Otherwise, push the current page on the history as marked, and move
 to physical page \arg.}
\end{tabularx}

\vspace{\stretch{1}}
\subsection*{Switching views}

\noindent
\begin{tabularx}{\linewidth}{clcX}
\ikey{w}{switch}{Switch view between master and client (if any)}
\ikey{W}{sync}{Goto page of client view corresponding to page of master view}
\ikey{\char94W}{autoswitch}{Toggle autoswitch flag}
\end{tabularx}

\vspace{\stretch{1}}

\subsection*{Table of contents}

\noindent
\begin{tabularx}{\linewidth}{clcX}
\ikey{T}{Thumbnails}{Process thumbnails.}
\ikey{t}{toc}{Display thumbnails if processed, or floating table of contents
 if available, or do nothing.}
\end{tabularx}

\vspace{\stretch{1}}

\vspace{\stretch{1}}

\newpage

\subsection*{Moving to pauses}

\noindent
\begin{tabularx}{\linewidth}{clcX}
\ikey{N}{next pause}{Move \arg pauses forward (equivalent to continue).}
\ikey{P}{previous pause}{Move \arg pauses backward.}
\end{tabularx}

\vspace{\stretch{1}}

\subsection*{Adjusting the page size}

\noindent
\begin{tabularx}{\linewidth}{clcX}
\ikey{\char94F}{fullscreen}{Adjust the size of the page to fit the
entire screen or reset the page to the default size (this is a toggle).}
\ikey{$<$}{smaller}{Scale down the resolution by scalestep (default
\tiny{$\sqrt{\sqrt{\sqrt 2}}$}).}
\ikey{$>$}{bigger}{Scale up the resolution by scalestep (default
\tiny{$\sqrt{\sqrt{\sqrt 2}}$}).}
\ikey{\char35}{fullpage}{Remove margins around the page and change
the resolution accordingly.}
\ikey{c}{center}{Center the page in the window, and resets the default
resolution.}
\end{tabularx}

\vspace{\stretch{1}}

\subsection*{Redisplay commands}

\noindent
\begin{tabularx}{\linewidth}{clcX}
\ikey{r}{redraw}{Redraw the current page to the current pause.}
\ikey{R}{reload}{Reload the file and redraw the current page.}
\ikey{\char94L}{redisplay}{Redisplay the current page to the first
pause of the page.}
\ikey{a}{active/passive}{toggle advi effects (so that reloading is silent).}
\end{tabularx}

\vspace{\stretch{1}}

\newpage

\vspace{\stretch{1}}

\subsection*{Moving the page in the window}

\noindent
\begin{tabularx}{\linewidth}{clcX}
\ikey{h}{page left}
{Moves one screen width toward the left of the page. Does nothing if the
  left part of the page is already displayed}
\ikey{l}{page right}
{Moves one screen width toward the right of the page. Does nothing if the
  right part of the page is already displayed}
\ikey{j}{page down}
{Moves one screen height toward the bottom of the page. Jumps to the top of
  next page, if there is one, and if the bottom of the page is already
  displayed.}
\ikey{k}{page up}
{Moves one screen height toward the top of the page. Jumps to the bottom
 previous page, if there is one, and if the top of the page is already
 displayed.}
\ikey{\char94 left button}{move page}
{A black line draws the page borders; moving the mouse then moves the
page in the window.}
\end{tabularx}

\vspace{\stretch{1}}

\newpage

\vspace{\stretch{1}}

\subsection*{Using the navigation history stack}

\noindent
\begin{tabularx}{\linewidth}{clcX}
\ikey{return}{forward}
{Push the current page on the history stack, and move forward n physical pages.}
\ikey{tab}{mark and next}
{Push the current page on the history as marked, and move forward n
physical pages.}
\ikey{backspace}{back}
{Move \arg pages backward according to the history. The history stack
is poped, accordingly.}
\ikey{escape}{find mark}
{Move \arg marked pages backward according to the history.
 Do nothing if the history does no contain any marked page.}
\end{tabularx}

\vspace{\stretch{1}}

\subsection*{Writing and drawing on the slide}

\noindent
\begin{tabularx}{\linewidth}{clcX}
\ikey{s}{scratch}{Give a pencil to type characters on the page.}
\ikey{S}{scratch}{Give a spray can to draw on the page.}
\end{tabularx}

\subsection*{Dealing with caches}

\noindent
\begin{tabularx}{\linewidth}{clcX}
\ikey{f}{load fonts}{Load all the fonts used in the document. By
default, fonts are loaded only when needed.}
\ikey{F}{make fonts}
{Does the same as \key{f}, and precomputes the glyphs of all the
characters used in the document.
This takes more time than loading the fonts, but the pages are drawn faster.}
\ikey{C}{clear}{Erase the image cache.}
\end{tabularx}

\vspace{\stretch{1}}






\def \newpage{\Newpage}

\section {Index}

\printindex

\end{document}

% LocalWords:  DVI Didier Weis advi INRIA LGPL Miquel Furuse Botlan XBuffer pst
% LocalWords:  XSelection Hyperref hyperref dvips advirecording advianchoring
% LocalWords:  CamlImages topleft topright bottomleft bottomright hbox PStricks
% LocalWords:  advianchor eepic pstricks ps pos bg colspec bubbleset inlined
% LocalWords:  flushpage XFlush XSync
